\section{Sharp constructions with arbitrary ramification}
\label{sec:constructions}

The following results answer a concrete existence question:
can complete cancellation of the perturbative cusp contribution
coexist with arbitrarily high inverse ramification?
The answer is yes, with a minimal number of factors.
The minimality statements are relative to finite products on
distinct positive dilations with nonzero integer exponents.

\subsection{Three eta factors realize every order}

\begin{theorem}[Minimal three-factor eta construction]
\label{thm:three-factors}
For each integer \(m\ge2\), put
\begin{align}
 K_m&=\frac{(m+1)^{(m^2-1)/2}}{m^{m^2/2}},
 &S_m&=\frac{4\pi^2}{m(m+1)},\notag\\
 F_m(t)&=
 K_m\,\frac{\eta(\ii t/(2\pi))
              \eta(\ii(m+1)t/(2\pi))^{m^2-1}}
             {\eta(\ii m t/(2\pi))^{m^2}}.
 \label{eq:eta-family}
\end{align}
For \(\Re t>0\),
\begin{equation}\label{eq:eta-family-dual}
 F_m(t)=R_m(Q),\qquad Q=\ee^{-S_m/t},\qquad
 R_m(Q)=
 \frac{P(Q^{m(m+1)})P(Q^m)^{m^2-1}}
      {P(Q^{m+1})^{m^2}}.
\end{equation}
The expansion begins
\begin{equation}\label{eq:eta-family-leading}
 R_m(Q)=1-(m^2-1)Q^m+m^2Q^{m+1}
                 +\bigO(Q^{m+2}).
\end{equation}
Thus the nome map has ramification order exactly \(m\), and its
support is primitive.  The same order occurs at every rational cusp
whose denominator is coprime to \(m(m+1)\), after the explicit
constant and phase normalization of \cref{eq:eta-cusp}.
Three distinct eta factors are necessary for a nonzero exponent
vector satisfying both
\(\sum r_d=0\) and \(\sum r_d/d=0\).
\end{theorem}
\begin{proof}
On the dilations \(1,m,m+1\), the exponent vector is
\[
 (r_1,r_m,r_{m+1})=(1,-m^2,m^2-1).
\]
Its two balances are
\[
 1-m^2+(m^2-1)=0,\qquad
 1-m+\frac{m^2-1}{m+1}=0.
\]
Applying \cref{eq:eta-cusp} at \(h=0,k=1\) cancels the
power and essential exponential.  The remaining constant is
\(K_m^{-1}\), proving \cref{eq:eta-family-dual}.
The factor with smallest dual exponent is \(P(Q^m)^{m^2-1}\);
the next is \(P(Q^{m+1})^{-m^2}\).
For \(m\ge2\), \(2m\ge m+2\), so no nonlinear product interferes
with their first two coefficients.  This proves
\cref{eq:eta-family-leading}.  Since \(m,m+1\) both occur,
their gcd is one.

If \((k,m(m+1))=1\), all \(g_d=1\), and the same two balances
hold.  The dual exponents remain \(m(m+1),m+1,m\).
The leading coefficients become
\(-(m^2-1)\rho_{m+1}\) and \(m^2\rho_m\).
Both are nonzero, so the order and primitive support persist.

With only two distinct dilations \(d_1,d_2\), the first balance
forces \(r_2=-r_1\), and the second gives
\(r_1(1/d_1-1/d_2)=0\).
Thus both exponents vanish.  One factor is even more immediate.
The three-factor construction attains the bound.
\end{proof}

For completeness, the unramified order \(1\) is realized by
\[
 \sqrt2\,\frac{\eta(\ii t/(2\pi))\eta(4\ii t/(2\pi))^2}
                  {\eta(2\ii t/(2\pi))^3}
 =\frac{P(Q^4)P(Q)^2}{P(Q^2)^3}
 =1-2Q+\bigO(Q^2),\qquad Q=\ee^{-\pi^2/t}.
\]
The exponent vector \((1,-3,2)\) on dilations \((1,2,4)\)
satisfies the same two balances.

There is no assertion here that every member of
\cref{eq:eta-family} has trivial modular character on
\(\Gamma_0(m(m+1))\).
In fact
\[
 \sum d r_d=m(m-1)
\]
need not be divisible by \(24\).
The eta multiplier transformation is sufficient for the proof;
the functions are analytic eta quotients with the stated
multipliers.

For the normalized output
\[
 Y=\frac{1-F_m(t)}{m^2-1},\qquad u^m=Y,
\]
the first inverse coefficient is
\[
 Q=u+\frac{m}{m^2-1}u^2+\bigO(u^3).
\]
Every subsequent coefficient is obtained by
\cref{eq:psi-lagrange}.
The principal radial branch has
\[
 t\sim \frac{mS_m}{\log(1/Y)}
       =\frac{4\pi^2/(m+1)}{\log(1/Y)}.
\]
Thus an unbounded finite ramification order coexists with a
logarithmic rather than an algebraic leading inverse in \(t\).

\subsection{The smallest nontrivial eta example}

For \(m=2\),
\[
 F_2(t)=\frac{3\sqrt3}{4}\,
 \frac{\eta(\ii t/(2\pi))\eta(3\ii t/(2\pi))^3}
      {\eta(2\ii t/(2\pi))^4},
 \qquad Q=\ee^{-2\pi^2/(3t)}.
\]
The exact dual series is
\begin{align}
 R_2(Q)={}&1-3Q^2+4Q^3-12Q^5+18Q^6
             -39Q^8+56Q^9\notag\\
          &\quad{}-108Q^{11}+148Q^{12}+\bigO(Q^{14}).
 \label{eq:R2}
\end{align}
With \(s=((1-F_2)/3)^{1/2}\), the branch \(Q\sim s\) is
\begin{align}
 Q={}&s+\frac23s^2+\frac{10}{9}s^3+\frac{10}{27}s^4
       -\frac{17}{27}s^5-\frac{1168}{243}s^6\notag\\
    &-\frac{4663}{486}s^7-\frac{11086}{729}s^8
       -\frac{48499}{13122}s^9
       +\frac{879718}{19683}s^{10}
       +\bigO(s^{11}).
 \label{eq:eta-example-Q}
\end{align}
The corresponding logarithmic correction is
\begin{align}
 \ell(s)=\log(Q/s)={}&
 \frac23s+\frac89s^2-\frac{22}{81}s^3-\frac{85}{81}s^4
 -\frac{1666}{405}s^5-\frac{25199}{4374}s^6
 +\bigO(s^7).\label{eq:eta-example-ell}
\end{align}
Consequently
\[
 t=\frac{2\pi^2/3}{-\log s-\ell(s)}.
\]
Replacing \(s\) by \(-s\) gives the second nome branch; the
logarithm lifts supply the remaining sheets.
All displayed rational coefficients were also checked by direct
composition, independently of the analytic proof.

The type depends on the cusp.  For the same normalized eta
quotient, at
\(\tau=1/2+\ii t/(2\pi)\), the exact identity is
\begin{equation}\label{eq:eta-minus-one}
 \frac{3\sqrt3}{4}
 \frac{\eta(\tau)\eta(3\tau)^3}{\eta(2\tau)^4}
 =
 \frac{\ee^{\pi\ii/12}}4\,
 \ee^{\pi^2/(4t)}
 \frac{P(-Q^6)P(-Q^2)^3}{P(Q^{12})^4},
 \qquad Q=\ee^{-\pi^2/(6t)}.
\end{equation}
The dual factor begins \(1+3Q^2+\cdots\).
Hence the flat finite cusp at \(q=1\) becomes an exponentially
growing cusp at \(q=-1\).
This is a direct illustration of why the arithmetic factors
\(g_d\) in the normal form matter.

\subsection{Four raw Euler-product factors: all perturbative terms cancel}

For an unnormalized Euler product, two balances leave an
\(\ee^{Dt/24}\) factor.  The following construction removes that
term too, without inserting a compensating exponential by hand.

\begin{theorem}[Minimal four-factor construction at every order]
\label{thm:four-factors}
Let \(m\ge1\),
\[
 g=\gcd(m,3),\qquad
 c_j=m+j,\qquad
 e_j=\frac{(-1)^j\binom3j(m+j)}{g}\quad (0\le j\le3),
\]
and put
\[
 L=\lcm(m,m+1,m+2,m+3),\qquad
 d_j=\frac{L}{c_j},\qquad
 C_m=\prod_{j=0}^3c_j^{e_j/2},\qquad
 T_m=\frac{4\pi^2}{L}.
\]
The exponents \(e_j\) are primitive nonzero integers, and
the raw product
\begin{equation}\label{eq:raw-family}
 E_m(t)=\prod_{j=0}^3P(\ee^{-d_jt})^{e_j}
\end{equation}
has the exact form
\begin{equation}\label{eq:raw-family-dual}
 E_m(t)=C_m\,\mathcal R_m(\ee^{-T_m/t}),\qquad
 \mathcal R_m(Q)=\prod_{j=0}^3P(Q^{c_j})^{e_j}
              =1-\frac{m}{g}Q^m+\bigO(Q^{m+1}).
\end{equation}
All three original balances vanish:
\[
 \sum e_j=0,\qquad
 \sum d_je_j=0,\qquad
 \sum e_j/d_j=0.
\]
The nome multiplicity is exactly \(m\).
Four distinct factors are necessary for a nonzero exponent
vector satisfying all three balances.
For \(m\ge2\), the next coefficient is
\(3(m+1)/g\), so the effective nome support is primitive.
For \(m=1\), primitiveness follows already from the nonzero
linear coefficient.
\end{theorem}
\begin{proof}
The third finite difference of a polynomial of degree at most
two vanishes.  Applying it to \(1,x,x^2\) gives
\[
 \sum_j\frac{e_j}{c_j}=0,\qquad
 \sum_je_j=0,\qquad
 \sum_jc_je_j=0.
\]
The gcd before dividing by \(g\) divides both \(m\) and \(m+3\),
and is therefore a divisor of \(\gcd(m,3)\).
Conversely \(\gcd(m,3)\) divides all four entries, so the
normalized vector is primitive.

The three balances in the original dilations follow from
\(d_j=L/c_j\):
\[
 \sum d_je_j=L\sum e_j/c_j=0,\qquad
 \sum e_j/d_j=L^{-1}\sum c_je_j=0.
\]
Apply \cref{eq:raw-cusp} at \(q=1\).
Its constant is
\[
 \prod_jd_j^{-e_j/2}
 =L^{-\sum e_j/2}\prod_jc_j^{e_j/2}=C_m,
\]
and its dual powers are \(Q^{c_j}\).
Only the \(j=0\) factor contributes in degree \(m\),
giving \(-m/g\).  For \(m\ge2\), the next degree
\(m+1\) has only the \(j=1\) contribution, equal to
\(3(m+1)/g\).
For \(m=1\), overlapping contributions change that coefficient
to \(5\), but do not change the leading order.

To prove minimality, consider at most three distinct positive
supports \(x_i\).  The balance matrix has rows
\(1,x_i,1/x_i\).
Multiplying its \(i\)-th column by \(x_i\) gives rows
\(x_i,x_i^2,1\), a Vandermonde matrix up to row order.
Its columns are independent for at most three distinct
nonzero \(x_i\).  Hence its nullspace is zero.
Our four nonzero factors attain the lower bound.
\end{proof}

The same proof classifies all rational balanced vectors on four
distinct positive rational supports \(c_0,c_1,c_2,c_3\):
they are proportional to
\begin{equation}\label{eq:four-support-classification}
 e_j=\frac{c_j}{\prod_{\ell\ne j}(c_j-c_\ell)}.
\end{equation}
Indeed Lagrange interpolation gives
\(\sum_jc_j^r/\prod_{\ell\ne j}(c_j-c_\ell)=0\)
for \(r=0,1,2\).
These are exactly the three balances after multiplying by the
appropriate power of \(c_j\).
The balance matrix has rank three, so the nullspace is
one-dimensional.  Clearing denominators and dividing by the
gcd gives its primitive integer vector.

For \(m=3\), the construction yields dual data
\[
 (c_j)=(3,4,5,6),\quad (e_j)=(1,-4,5,-2),\quad
 (d_j)=(20,15,12,10).
\]
Here \(L=60\), \(T_3=\pi^2/15\), and
\(C_3=25\sqrt{15}/96\).
The original four raw products have a constant perturbative
expansion and a cubic nome defect.

For every \(m\ge2\), writing
\[
 y=\frac{1-E_m/C_m}{m/g},\qquad s^m=y,
\]
gives
\[
 Q=s+\frac{3(m+1)}{m^2}s^2+\bigO(s^3).
\]
The complete inverse and its certified tail are
\cref{eq:all-flat-sheets,eq:flat-certified}.
The construction therefore supplies a family of explicit
\(q\)-analog critical points of every order, together with
convergent inverse expansions.

\subsection{A low-level raw product}

A compact additional example is
\[
 E(t)=
 \frac{P(\ee^{-2t})^5P(\ee^{-6t})}
      {P(\ee^{-t})P(\ee^{-3t})^5}.
\]
The exponent vector on \(1,2,3,6\) is \((-1,5,-5,1)\);
its three balances are zero.  Hence
\[
 E(t)=\frac98
 \frac{P(Q)P(Q^3)^5}{P(Q^2)^5P(Q^6)},
 \qquad Q=\ee^{-2\pi^2/(3t)}.
\]
The normalized dual product is
\(1-Q+4Q^2-10Q^3+20Q^4-39Q^5+\cdots\).
It gives a simple nome inverse and a logarithmic \(t\)-inverse.
This example needs no eta prefactor and makes the distinction
between raw and normalized observables concrete.
